% Style based on the supplied Waseem_Generic_AI_ML_GenAI_AgenticAI.tex.
% Python replaces template tokens. Edit layout here, content in JSON/portfolio JS.
\documentclass[9pt,letterpaper]{extarticle}
\usepackage[left=0.68in,right=0.68in,top=0.45in,bottom=0.45in]{geometry}
\usepackage[T1]{fontenc}
\input{glyphtounicode}
\pdfgentounicode=1
\usepackage[utf8]{inputenc}
\usepackage{lmodern}
\usepackage{microtype}
\usepackage{pifont}
\usepackage{xcolor}
\usepackage{enumitem}
\usepackage{titlesec}
\usepackage[hidelinks]{hyperref}
\usepackage{tabularx,array,parskip,ragged2e,needspace}
\pagestyle{empty}
\setlength{\parindent}{0pt}
\setlength{\parskip}{1pt}
\definecolor{accent}{HTML}{155B8A}
\titleformat{\section}{\large\bfseries\color{accent}}{}{0pt}{}[\vspace{0.5pt}\color{accent}\titlerule]
\titlespacing*{\section}{0pt}{3pt}{2pt}
% Compact subsections beneath the main accent-rule sections.
\titleformat{\subsection}{\normalsize\bfseries\color{accent}}{}{0pt}{}
\titlespacing*{\subsection}{0pt}{3pt}{1pt}
\setlist[itemize]{leftmargin=1.15em,itemsep=0.8pt,topsep=1pt,parsep=0pt,partopsep=0pt}
\newcommand{\role}[4]{%
  \Needspace{4\baselineskip}\textbf{#1}\hfill\textbf{#2}\\[-1pt]
  \textit{#3}\hfill\textit{#4}\par}
\newcommand{\entry}[2]{\Needspace{4\baselineskip}\textbf{#1}\hfill\textit{#2}\par}
\newcommand{\skillrow}[2]{\textbf{#1:} & #2\\}
\hypersetup{pdftitle={Waseem Raza, PhD - Comprehensive 3-Page Resume},pdfauthor={Waseem Raza, PhD}}
\begin{document}

\begin{center}
{\LARGE\bfseries Waseem Raza, PhD}{\normalsize\mdseries }\\[2pt]
{\large\bfseries Machine Learning Engineer \textbar{} AI/ML, GenAI, Data \& Intelligent Systems}\\[1pt]
{\small Applied AI \textbar{} Generative AI \textbar{} Wireless Networks \textbar{} Research and Engineering}\\[3pt]
{\footnotesize \ding{41}\enspace \href{mailto:waseem.placements@gmail.com}{waseem.placements@gmail.com} \textbar{} \href{mailto:waseemraza844@gmail.com}{waseemraza844@gmail.com}\\[1pt]New Jersey, United States \quad \ding{37}\enspace \href{tel:+17328039703}{+1 732-803-9703}\\[1pt]\href{https://www.linkedin.com/in/waseem-raza-phd-9a737858/}{LinkedIn} \textbar{} \href{https://github.com/WaseemRaza844}{GitHub} \textbar{} \href{https://scholar.google.com/citations?user=d58Drr0AAAAJ\&hl=en}{Google Scholar} \textbar{} \href{https://waseemraza844.github.io/portfolio/learning.html}{Coursera Courses} \textbar{} \href{https://waseemraza844.github.io/portfolio/}{Portfolio}}
\end{center}
\vspace{-5pt}
\Needspace{3\baselineskip}
\section{Professional Profile}
AI/ML and Generative AI Engineer with a PhD in Electrical and Computer Engineering and 7+ years of combined industry, research, and engineering experience across enterprise AI, data science, wireless systems, and applied research. Experience spans LLM-enabled analytics, retrieval-augmented generation (RAG), domain-adaptive machine learning, predictive modeling, time-series analysis, ETL/ELT, Azure Databricks, Apache Spark, MLOps, 5G/6G network intelligence, and experimental research. At AT\&T, lead development and operational integration of AI solutions for critical service-assurance workflows, combining model validation, data quality, cross-domain evaluation, and stakeholder-driven delivery.
\Needspace{2\baselineskip}
\subsection{Technical Competencies}
\Needspace{2\baselineskip}\noindent{\small\textbf{Generative \& Agentic AI:} LLM applications, RAG, embeddings, semantic retrieval, prompt engineering, structured extraction, LangChain, LangGraph, LlamaIndex, tool-using agents, memory, orchestration, and grounded generation}\par
\Needspace{2\baselineskip}\noindent{\small\textbf{Machine Learning:} PyTorch, TensorFlow, scikit-learn, XGBoost, classification, regression, anomaly detection, time-series forecasting, transfer learning, domain adaptation, cross-validation, and model evaluation}\par
\Needspace{2\baselineskip}\noindent{\small\textbf{Data Engineering \& Analytics:} Python, SQL, Pandas, NumPy, ETL/ELT, feature engineering, data validation, Azure Databricks, Apache Spark, PySpark, Synapse Analytics, and cloud analytics}\par
\Needspace{2\baselineskip}\noindent{\small\textbf{MLOps \& Software Engineering:} Docker, Kubernetes, OpenShift, Git, GitHub Actions, CI/CD, REST APIs, Flask, testing, Linux, model deployment, system design, and reproducible workflows}\par
\Needspace{2\baselineskip}\noindent{\small\textbf{Cloud \& AI Platforms:} AWS Bedrock, Azure OpenAI, Microsoft Foundry, Azure AI Services, Hugging Face Transformers, Google Cloud analytics, managed AI services, and foundation-model workflows}\par
\Needspace{2\baselineskip}\noindent{\small\textbf{Wireless \& Network Intelligence:} LTE/4G, 5G-NR, NR-NTN, RAN, mobility, propagation modeling, network measurements, outage management, optimization, self-healing, and service assurance}\par
\Needspace{2\baselineskip}\noindent{\small\textbf{Research \& Experimentation:} 3GPP-aligned simulation, synthetic data, GANs, reinforcement learning, robustness testing, distribution-shift evaluation, statistical validation, and experimental design}\par
\Needspace{4\baselineskip}\textbf{Leadership and Professional Strengths}\par
\begin{itemize}
\item \textbf{Ownership and accountability:} Lead AI/ML initiatives through requirements analysis, data preparation, model development, validation, stakeholder review, and deployment planning.
\item \textbf{Technical leadership:} Define technical approaches, evaluate alternatives, coordinate experimentation, and communicate findings to guide projects from ambiguous requirements to measurable outcomes.
\item \textbf{Cross-functional collaboration:} Work with software engineers, network specialists, researchers, geologists, and public-safety stakeholders to translate domain requirements into practical data and AI solutions.
\item \textbf{Continuous learning:} Apply ongoing study in generative AI, agentic AI, RAG, LLM evaluation, and cloud-native platforms through hands-on projects and professional coursework.
\item \textbf{Reliability and commitment:} Maintain reproducible workflows, careful validation, and attention to detail throughout experimentation and delivery.
\item \textbf{Adaptability and problem solving:} Combine critical thinking, systematic experimentation, and data-driven analysis to address unfamiliar technologies and complex operational problems.
\end{itemize}
\Needspace{3\baselineskip}
\section{Professional Experience}
\role{Machine Learning Engineer}{02/2025 - Present}{AT\&T, contractor via Innovatix Inc.}{Bedminster, NJ}
\begin{itemize}\interlinepenalty=10000
\item Own end-to-end development, testing, and operational integration of an AI-driven diagnostic system for critical E911 service-assurance workflows, improving automated root-cause classification accuracy from 20\% to 80\%.
\item Designed domain-adaptive neural models and reusable data-processing workflows that generalize across heterogeneous operational domains and reduce performance degradation caused by distribution shifts.
\item Developed multimodal AI prototypes combining work-order text, location information, network context, and operational signals to support GenAI-assisted diagnostics for VoNR and 5G service assurance.
\item Created model-validation, error-analysis, and failure-mode assessment procedures to compare performance across domains and support reliable deployment in high-impact operational environments.
\item Partnered with network engineers, software teams, and public-safety stakeholders to translate operational requirements into validated AI workflows and production-oriented data assets.
\end{itemize}
\role{Data Scientist and AI/ML Engineer}{09/2024 - 12/2024}{Impac Exploration Services}{Weatherford, OK}
\begin{itemize}\interlinepenalty=10000
\item Built an LLM-enabled analytical pipeline for extracting and interpreting information from geochemical well data.
\item Developed scalable ingestion, cleansing, transformation, and validation workflows using Python, Azure Databricks, Azure Functions, and Synapse Analytics.
\item Automated ETL processing for large exploration datasets, reducing data-preparation time by 40\% and improving reproducibility.
\item Implemented cross-validation, data-quality, and consistency checks that improved reliability of analytical outputs by 20\%.
\item Collaborated with geologists to convert domain requirements and technical information into AI-ready datasets and actionable analytical outputs.
\end{itemize}
\role{Advanced-Degrees Intern}{06/2024 - 08/2024}{AT\&T Labs}{Austin, TX}
\begin{itemize}\interlinepenalty=10000
\item Analyzed 500K+ LTE/5G-NR diagnostic logs and applied ML methods to detect anomalies, reducing diagnostic time by 20\%.
\item Investigated 5G NR-NTN cell-selection behavior and improved decision algorithms by 15\% through data-driven analysis and validation.
\item Conducted structured analysis of large operational datasets and communicated technical findings to engineering teams supporting next-generation wireless systems.
\end{itemize}
\role{Advanced Wireless Technology Development Intern}{09/2022 - 01/2023}{Futurewei Technologies}{Schaumburg, IL}
\begin{itemize}\interlinepenalty=10000
\item Improved autoencoder-based channel-state-information compression by 25\% using transfer learning and systematic experimentation on large wireless datasets.
\item Conducted 50+ experiments to validate model performance, robustness, and generalization under changing data conditions.
\item Collaborated with research and engineering stakeholders to translate modeling results into recommendations for next-generation communication systems.
\end{itemize}
\role{Advanced-Degrees Intern}{06/2022 - 08/2022}{AT\&T Labs}{Bedminster, NJ}
\begin{itemize}\interlinepenalty=10000
\item Designed ML and statistical methods for analyzing users' timing in 5G standalone networks, improving analytical precision by 20\%.
\item Analyzed large-scale user trace data with SQL and visualization tools and communicated findings to senior technical leadership.
\item Converted raw operational traces into interpretable performance indicators to support diagnosis and network optimization.
\end{itemize}
\role{Graduate Research Assistant}{09/2019 - 07/2025}{AI4Networks, University of Oklahoma}{Tulsa, OK}
\begin{itemize}\interlinepenalty=10000
\item Designed reusable Python simulation and data-generation platforms for large-scale experimentation, including a 3GPP-compliant network simulator and an AI-enabled 5G testbed.
\item Built transfer-learning, adversarial-learning, GAN, autoencoder, and reinforcement-learning pipelines for sparse-data and cross-domain modeling.
\item Developed robustness evaluation frameworks for models operating under distribution shifts, missing data, and limited labeled samples; led 100+ experiments from problem formulation through validation.
\item Developed scalable self-healing and outage-management frameworks that improved modeled network resilience by 30\%.
\item Published peer-reviewed research, presented technical results internationally, and mentored junior researchers in ML implementation, experimental design, and technical writing.
\end{itemize}
\Needspace{3\baselineskip}
\section{Education}
\entry{PhD, Electrical and Computer Engineering}{09/2019 - 07/2025}
University of Oklahoma, Norman, Oklahoma, USA\par
\entry{MSc, Telecommunication Engineering}{08/2014 - 08/2016}
University of Engineering and Technology, Taxila, Pakistan\par
\entry{BSc, Telecommunication Engineering}{11/2010 - 06/2014}
University of Engineering and Technology, Taxila, Pakistan\par
\Needspace{3\baselineskip}
\section{Selected Projects \& Research}
\Needspace{2\baselineskip}
\subsection{Projects \& Applied Research}
\entry{\href{https://github.com/WaseemRaza844/agentic-ai-projects-by-was}{Agentic AI and RAG Applications}}{Portfolio Project \textbar{} 2026}
Built grounded GenAI workflows that combine document retrieval, vector search, stateful orchestration, tool-using agents, memory, evaluation, and multi-agent handoffs. Emphasized controllable execution, traceable context flow, and reliable responses for multi-step tasks rather than single unconstrained model calls.\par
{\small\textit{Tools \& Methods:} Python, LangGraph, LangChain, CrewAI, AutoGen, RAG, Vector Databases, Embeddings, Tool Calling, Agent Memory, Evaluation\par}
\vspace{1pt}
\entry{AI-Ready Lakehouse Pipeline}{Portfolio Project \textbar{} 2024}
Designed an AI-ready lakehouse pipeline that prepares structured and semi-structured data for analytics and downstream machine-learning workloads. Implemented repeatable ingestion, cleansing, transformation, validation, and scalable processing patterns suitable for reproducible cloud data workflows.\par
{\small\textit{Tools \& Methods:} Azure Databricks, Apache Spark, PySpark, ETL/ELT, Data Ingestion, Schema Transformation, Data Validation, Feature Preparation, Cloud Data Engineering\par}
\vspace{1pt}
\entry{Financial ML and Forecasting}{Portfolio Project \textbar{} 2024-2025}
Developed finance-oriented machine-learning workflows for credit-risk and loan-default modeling together with time-series forecasting experiments. Applied data preparation, feature engineering, model comparison, cross-validation, and predictive-performance analysis to structured financial datasets.\par
{\small\textit{Tools \& Methods:} Python, scikit-learn, Classification, Credit Risk, Loan Default Modeling, Time-Series Forecasting, Feature Engineering, Cross-Validation, Model Evaluation\par}
\vspace{1pt}
\entry{Wearable Intelligence for Health Monitoring}{Research Project \textbar{} 2024-2026}
Developed deep-learning workflows for recognizing medication-taking and dietary activities from smartwatch accelerometer and gyroscope signals in free-living settings. Evaluated windowing, feature extraction, CNN/LSTM-style sequence models, cross-user generalization, and recall-oriented performance for health-monitoring applications.\par
{\small\textit{Tools \& Methods:} Wearable AI, Smartwatch IMU, Time Series, Deep Learning, CNN/LSTM, Feature Engineering, Cross-User Evaluation, Activity Recognition, Healthcare AI\par}
\vspace{1pt}
\entry{Propagation Models \& Digital Twins}{Research Project \textbar{} 2024}
Developed robust data-driven radio-propagation models for wireless digital twins using domain-informed generative learning and scarce measured data. Combined 3GPP-informed pretraining, GAN-based augmentation, domain adaptation, and statistical validation to improve resilience under distribution shift.\par
{\small\textit{Tools \& Methods:} Propagation Modeling, Digital Twins, GANs, Domain Adaptation, 3GPP Models, Synthetic Data, Statistical Validation, Wireless AI, Distribution Shift\par}
\vspace{1pt}
\entry{SyntheticNET: 3GPP-Compliant Network Simulation}{Research Project \textbar{} 2020-2024}
Developed a Python-based, 3GPP-aligned simulation environment for realistic mobility, propagation, and cellular-network experimentation when measured data is sparse or costly. Supported controlled data generation and repeatable studies in propagation modeling, network optimization, digital twins, and resilient wireless AI.\par
{\small\textit{Tools \& Methods:} Python, 3GPP, Wireless Simulation, Mobility Modeling, Propagation Modeling, Synthetic Data, Experiment Design, Digital Twins, Network Analytics\par}
\vspace{1pt}
\entry{\href{https://github.com/WaseemRaza844/system-design-masterclass}{System Design Case-Study Portfolio}}{Guided Project \textbar{} 2026}
Developed architecture case studies that translate ambiguous product requirements and traffic assumptions into service boundaries, APIs, data models, storage choices, caching strategies, and reliability targets. Documented scalability, consistency, capacity, and failure-mode tradeoffs to support production-oriented system design decisions.\par
{\small\textit{Tools \& Methods:} System Design, REST APIs, Distributed Systems, Data Modeling, Caching, Scalability, Reliability, Capacity Planning, Architecture Tradeoffs\par}
\vspace{1pt}
\Needspace{2\baselineskip}
\subsection{Publications}
\entry{Real-World Medication Adherence Monitoring Using Hybrid Neural Networks and Smartwatches}{2026}
\textit{HealthINF \textbar{} Accepted}\par
\entry{\href{https://doi.org/10.3390/s25247651}{On Multi-Parameter Optimization and Proactive Reliability in 5G and Beyond Cellular Networks}}{2025}
\textit{Sensors 2025}\par
\entry{\href{https://doi.org/10.1109/PIMRC59610.2024.10817403}{An AI-Driven Framework for Enhancing Resilience in Propagation Models to Enable Digital Twin}}{2024}
\textit{IEEE PIMRC}\par
\entry{Data-Driven Intelligent Outage Management for High Shadowing Environments in 5G\&B Networks}{2024}
\textit{IEEE SmartNets}\par
\entry{\href{https://arxiv.org/abs/2203.15899}{Toward a Hybrid RF/Optical Lunar Communication System (LunarComm)}}{2022}
\textit{IEEE Network}\par
\entry{Towards Positioning Error Impact Characterization and Minimization in User-Centric RAN}{2022}
\textit{IEEE WCNC}\par
\entry{\href{https://doi.org/10.1016/j.aeue.2020.153244}{Unification of RF Energy Harvesting Schemes under Mixed Rayleigh-Rician Fading Channels}}{2020}
\textit{AEÜ - International Journal of Electronics and Communications}\par
\Needspace{3\baselineskip}
\section{Certifications \& Professional Development}
\Needspace{3\baselineskip}\noindent\textbf{1. Generative AI Applications with Amazon Bedrock} {\small --- \textcolor{blue}{\underline{\href{https://www.coursera.org/account/accomplishments/verify/3W16P8EE4BJZ}{Coursera}}} \textbar{} Completed \textbar{} August 27, 2026 \textbar{} \textcolor{blue}{\underline{\href{https://waseemraza844.github.io/portfolio/certificates/genai/aws-bedrock/c2-generative-ai-applications-bedrock.pdf}{Certificate}}}}\par
\textbf{Focus:} Covers development of generative-AI applications on Amazon Bedrock, including foundation-model access, prompt-based application logic, model interaction, and managed AWS services. Emphasizes practical cloud patterns for building and operating LLM-powered applications.\par
{\small\textbf{Tools \& Skills:} Amazon Bedrock, foundation models, prompt engineering, AWS generative-AI services, LLM application development, cloud AI workflows\par}
\vspace{2pt}
\Needspace{3\baselineskip}\noindent\textbf{2. Agentic AI with LangChain and LangGraph} {\small --- \textcolor{blue}{\underline{\href{https://www.coursera.org/account/accomplishments/verify/L4EY4SLGCPUE}{Coursera}}} \textbar{} Completed \textbar{} February 9, 2026}\par
\textbf{Focus:} Covers design of stateful agentic applications using LangChain and LangGraph, including graph-based orchestration, tool use, memory and state, routing, and multi-step workflows. Emphasizes controllable agent execution and composition of reliable LLM-driven processes.\par
{\small\textbf{Tools \& Skills:} LangChain, LangGraph, tool-calling, state and memory, workflow orchestration, routing, agent design\par}
\vspace{2pt}
\Needspace{3\baselineskip}\noindent\textbf{3. Machine Learning} {\small --- \textcolor{blue}{\underline{\href{https://www.coursera.org/account/accomplishments/verify/VU3Y5TBLVLZT}{Stanford Online / Coursera}}} \textbar{} Completed \textbar{} March 29, 2020}\par
\textbf{Focus:} Covers foundational machine-learning methods for supervised and unsupervised problems, including regression, classification, model evaluation, regularization, clustering, and practical model-selection considerations. Emphasizes translating data into predictive models while controlling generalization error.\par
{\small\textbf{Tools \& Skills:} supervised learning, regression, classification, regularization, clustering, model evaluation, bias-variance tradeoffs, machine-learning workflows\par}
\vspace{2pt}
\Needspace{3\baselineskip}\noindent\textbf{4. Introduction to Data Analytics in Google Cloud} {\small --- \textcolor{blue}{\underline{\href{https://www.coursera.org/account/accomplishments/verify/R9FGPRZ3NJE8}{Coursera}}} \textbar{} Completed \textbar{} July 24, 2024}\par
\textbf{Focus:} Introduces cloud-native data analytics in Google Cloud, covering the analytics lifecycle, data sources, storage and processing concepts, and how cloud services support scalable analysis. Emphasizes organizing data work to produce reliable analytical insight.\par
{\small\textbf{Tools \& Skills:} Google Cloud, cloud analytics, data lifecycle, data-ingestion concepts, storage and processing, analytical workflows, data-driven decision support\par}
\vspace{2pt}
\Needspace{3\baselineskip}\noindent\textbf{5. Introduction to Deep Learning \& Neural Networks with Keras} {\small --- \textcolor{blue}{\underline{\href{https://www.coursera.org/account/accomplishments/verify/Y6NP4GWLS4DZ}{IBM / Coursera}}} \textbar{} Completed \textbar{} August 12, 2026}\par
\textbf{Focus:} Introduces deep-learning and neural-network fundamentals using Keras, covering network architecture, training, validation, and model improvement. Connects concepts such as activation functions and optimization with practical neural-network implementation and evaluation.\par
{\small\textbf{Tools \& Skills:} Keras, deep learning, neural networks, model training, validation, activation functions, optimization, model evaluation\par}
\vspace{2pt}
\end{document}
